\begin{pblock}{Cloud screening that keeps the data}
On an Atlantic coast a scene-level cloud filter is ruinous: discard every
image with clouds and little survives. But a cloud over a hillside three
kilometres inland does not obscure the flat.

So cloudiness is measured \term{only where it matters}. Two streaming passes:

\step{1}{every pixel votes across the archive — permanent water, permanent
land, or \term{transition}, which is the intertidal candidate;}
\step{2}{cloud cover is then computed \emph{inside the transition zone
alone}, with a dilated safety buffer.}

Villaviciosa, nine years, 1379 acquisitions: only 141 are
essentially clear over the flat. Demanding clear scenes would leave a decade
of imagery yielding almost nothing.

\vspace{2mm}
But cloud cover is not a property of a scene — it is a property of each pixel.
Masking pixel by pixel inside the transition zone lets partially clouded
acquisitions still contribute: \emphnum{415} scenes have at least half
the flat visible, and every clear pixel in them enters the fit. The estimator
needs a minimum number of observations \emph{per pixel}, not per scene, which
is what makes that possible.
\end{pblock}
